\section{Exact second moment for the dipole--composite probe}
\label{app:general-moment}

This appendix assumes a uniform spin-$s$ nearest-neighbor periodic chain
with even $N\ge4$, $s\ge1/2$, and the probe form specified before
Eq.~\eqref{eq:general-moment}. Set $c=s(s+1)$ and $d=2s+1$.
On one bond, the quadrupole--spin vector obeys
\begin{equation}
 X_{ij}=\tfrac12\{h_{ij},D_{ij}\}+\tfrac{c}{3}D_{ij},
 \qquad D_{ij}=T_i^z-T_j^z.
\end{equation}
In the $J-1\to J$ channel its coefficient relative to $D_{ij}$ is
$r_J=J^2/2-2c/3$. The squared $D^z$ matrix elements summed over magnetic
components are $J(d^2-J^2)/3$, and the corresponding energy difference is
$J$. Including both transition directions, define
$w_J=2J^3(d^2-J^2)/3$ for $1\le J\le d-1$. Thus the Gram entries of
$C_D=[h_{ij},D_{ij}]$ and $C_X=[h_{ij},X_{ij}]$ are sums of
$w_J$, $w_Jr_J$, and $w_Jr_J^2$:
\begin{align}
 G_{DD}&=\frac{d^2(d^2-1)^2}{18},\notag\\
 \frac{G_{DX}}{G_{DD}}&=\frac{d^2-4}{12},\qquad
 \frac{G_{XX}}{G_{DD}}=\frac{7(d^2-4)^2}{360}.
 \label{eq:general-bond-gram}
\end{align}
These finite sums follow from power-sum identities. We verified the
identities coefficient by coefficient through their finite differences,
rather than fitting values at selected $s$. The minimized bond norm is
$d^2(d^2-1)^2(d^2-4)^2/1440$.
For $s=3/2$, the Gram matrix is
\begin{equation}
 \begin{pmatrix}G_{DD}&G_{DX}\\G_{DX}&G_{XX}\end{pmatrix}
 =\begin{pmatrix}200&200\\200&560\end{pmatrix},
 \label{eq:thermal-gram}
\end{equation}
with minimized bond norm $360$.

The adjacent-bond contribution on sites $0,1,2$ is
\begin{equation}
 Z_{012}=[h_{12},X_{01}]-[h_{01},X_{12}].
\end{equation}
With repeated Cartesian indices summed, its explicit decomposition is
\begin{align}
 Z_{012}/i={}&\epsilon_{abc}
  (Q_0^{zb}T_1^cT_2^a+T_0^aT_1^cQ_2^{zb})\notag\\
 &-T_0^a(\epsilon_{bzc}Q_1^{ca}
                  +\epsilon_{azc}Q_1^{cb})T_2^b.
 \label{eq:three-site-commutator}
\end{align}
The three terms have local tensor-rank patterns $QTT$, $TTQ$, and
$TQT$, respectively. They are mutually orthogonal because
$\operatorname{tr}(T^aQ^{bc})=0$. Rotational invariance and the traces of
$(T^z)^2$ and $(T^z)^4$ give
\begin{align}
 \operatorname{tr}(T^aT^b)&=t\delta_{ab},\quad t=dc/3,\notag\\
 \operatorname{tr}(Q^{ab}Q^{ef})&=v(\delta_{ae}\delta_{bf}
 +\delta_{af}\delta_{be}-\tfrac23\delta_{ab}\delta_{ef}),\notag\\
 v&=dc(4c-3)/60.
\end{align}
Contracting the Levi-Civita symbols in
Eq.~\eqref{eq:three-site-commutator} gives $20/3$, $20/3$, and $20$
times $vt^2$ for those three patterns. Therefore
\begin{equation}
 \|Z_{012}\|_{\rm HS}^2=\frac{100}{3}vt^2
                      =\frac{5}{81}d^3c^3(4c-3).
 \label{eq:three-site-norm}
\end{equation}

All one-site partial traces of $C_D,C_X$ and all two-site partial traces
of $Z$ vanish. Distinct two-site and three-site supports are consequently
orthogonal in the full uniform trace. For a periodic chain with $N\ge4$,
there are $N$ distinct supports of each type. The alternating signs give
\begin{equation}
 M_2(u)=\frac{\|bC_X+uC_D\|_{\rm HS}^2}{\ell d^2}
       +\frac{b^2\|Z\|_{\rm HS}^2}{\ell d^3}.
\end{equation}
Completing the square with Eqs.~\eqref{eq:general-bond-gram} and
\eqref{eq:three-site-norm} yields Eq.~\eqref{eq:general-moment}.
The first contribution alone is the earlier lower bound. The second
cannot be tuned away by changing $u$.
Keeping only the two-site contribution already gives the quartet bound
\begin{equation}
 \min_u\frac{\|[h,R]\|_{\rm HS}^2}{4^N}
 \ge\frac{360b^2}{5\times16}=\frac92b^2,
 \label{eq:thermal-bound}
\end{equation}
which the exact result sharpens to $(197/16)b^2$.

The local identities and partial traces were additionally checked with
exact rational matrices for $s=1/2,1,3/2,2,5/2,3$. Separate orthonormal
Cartesian-spin Kronecker matrices on the full $N=4,6$ chains, without
importing the rational proof routines, verify the whole quadratic for
$s=1/2,1,3/2$ to floating-point accuracy. These checks support the
algebraic derivation; they do not extend its assumptions to other probe
structures, open chains, odd rings, or nonuniform exchange.

\subsection{A size-uniform positive-frequency weight bound}
\label{app:weight-bound}

Use the normalized Hilbert--Schmidt norm on each operator's support,
$\|O\|_2^2=d^{-|\operatorname{supp}O|}\operatorname{Tr}O^\dagger O$.
The full uniform-trace spectral measure of $R$ is symmetric in frequency.
Define its moments $M_{2k}$ including both frequency signs; $M_0$ also
includes equal-energy matrix elements. Local tensor orthogonality gives
\begin{equation}
 M_0=\frac{1}{\ell}\left[\frac{c}{3}u^2
          +\frac{c^2(4c-3)}{27}b^2\right].
 \label{eq:uniform-M0}
\end{equation}
For $s\ge1$, comparison with Eq.~\eqref{eq:general-moment} yields
$M_2\ge(3c/2)M_0$. To see positivity without optimizing an irrational
eigenvalue, rescale $b,u$ by the square roots of their coefficients in
$\ell M_0$. The matrix for $\ell[M_2-(3c/2)M_0]$ in these coordinates has
diagonal entries $(61c-42)/30$ and $7c/6$, and determinant
$c(107c-54)/180$. All are positive for $c\ge2$.

For the fourth moment write
\begin{align}
 [h,R]&=\frac{1}{\sqrt{\ell N}}\sum_i(-1)^i B_i,\notag\\
 B_i&=bC_{X,i}+uC_{D,i}+bZ_{i,i+1,i+2}.
\end{align}
The two-site and three-site terms are orthogonal as above, giving
$\|B_i\|_2^2=\ell M_2$. Only the four bonds meeting sites $i,i+1,i+2$
contribute to $G_i=[h,B_i]$. Since $\|h_{ij}\|_{\rm op}=c$,
$\|G_i\|_2\le8c\|B_i\|_2$. Each traceless $G_i$ is supported within
the five-site interval from $i-1$ to $i+3$. Two disjoint intervals have
zero uniform-trace inner product. At most nine translated intervals,
including itself, overlap a given interval; for $N=4,6,8$ the number is
even smaller. Bounding the remaining inner products by their norms gives
\begin{equation}
 M_4\le\frac{9}{\ell}\|G_i\|_2^2
      \le576c^2M_2.
 \label{eq:uniform-M4-bound}
\end{equation}
This is a locality bound, not a fit to numerical moments.

Separate the spectral measure into $|\omega|\le\nu$ and
$|\omega|>\nu$. The low-frequency contribution to $M_2$ is at most
$\nu^2M_0$. Cauchy--Schwarz for the remaining measure therefore gives
\begin{equation}
 C^{\rm tr}_{>\nu}\ge
 \frac{[M_2-\nu^2M_0]_+^2}{2M_4},
 \label{eq:window-moment-inequality}
\end{equation}
where $[x]_+=\max(x,0)$ and the factor two accounts for the two frequency
signs. Taking $\nu=\sqrt{3c/4}$ and using the preceding inequalities,
\begin{align}
 C^{\rm tr}_{>\sqrt{3c/4}}
 &\ge\frac{M_2}{4608c^2}\notag\\
 &\ge\frac{b^2(4c-3)(86c-27)}{3732480\ell}>0
 \quad(s\ge1,\ b\ne0).
 \label{eq:general-window-bound}
\end{align}
This holds uniformly over every real retuning $u$ and every finite even
$N\ge4$. It is the ordered uniform-trace limit of the effective model,
not an assertion about arbitrary finite scaled temperature or the full
microscopic infinite-temperature spectrum.

Full-chain matrices at $s=1,3/2$ and $N=4,6$ verify the zeroth-, second-,
and fourth-moment inequalities. Complete four-site diagonalization also
checks Eq.~\eqref{eq:window-moment-inequality} from actual transition
weights. Minimizing the quartet window weight gives $0.350008b^2$,
about 1842 times its conservative lower bound. Thus the lower bound is
useful as a size-uniform nonvanishing statement, not as a quantitative
estimate of the scattering intensity.
